\documentclass[conference]{IEEEtran}
\usepackage{cite}
\usepackage{amsmath,amssymb,amsfonts}
\usepackage{algorithmic}
\usepackage{graphicx}
\usepackage{textcomp}
\usepackage{xcolor}
\usepackage{booktabs}
\usepackage{url}
\usepackage{microtype}

\begin{document}

\title{Reduced-Scale Reproduction and Ablation Study of Barlow Twins on CIFAR-10}

\author{\IEEEauthorblockN{Khagesh Ranjan}
\IEEEauthorblockA{\textit{Department of Computer Science and Engineering} \\
\textit{Rajiv Gandhi Institute of Petroleum Technology (RGIPT)} \\
\textit{Roll Number: 24cs2021} \\
\textit{DL-070 Deep Learning Laboratory -- October 2026}}
}

\maketitle

\begin{abstract}
Self-supervised visual representation learning has traditionally relied on either large-minibatch contrastive objectives requiring thousands of negative sample pairs, or asymmetric architectures with heuristic stop-gradient mechanics. Barlow Twins introduces an alternative, biologically inspired paradigm: driving the cross-correlation matrix between twin view embeddings toward the identity matrix, simultaneously maximizing agreement across augmented views while minimizing redundancy across feature dimensions. In this work, we conduct a rigorous reduced-scale reproduction and ablation study of Barlow Twins on the CIFAR-10 dataset using consumer-grade local hardware (an NVIDIA GeForce RTX 5050 Laptop GPU). Under a constrained 50-epoch budget ($N=64, D=512$), standard Barlow Twins attains 74.84\% linear probe test accuracy, 72.41\% 20-NN accuracy, and an effective feature representation rank of 506/512. In our primary ablation study, removing the redundancy reduction objective entirely ($\lambda=0$) induces a severe 17.72 percentage point drop in linear probe accuracy (down to 57.12\%) and a 20.43 percentage point drop in $k$-NN accuracy, demonstrating empirically that invariance alone is incapable of supporting effective semantic representations. Furthermore, we uncover a catastrophic failure mode in static hinge thresholding, where zeroing out weak off-diagonal penalties causes full dimensional collapse to rank 0 and 10.00\% chance-level accuracy. Finally, we analyze minibatch estimation variance under small batches and outline promising directions for future exploration.
\end{abstract}

\begin{IEEEkeywords}
Self-Supervised Learning, Barlow Twins, Redundancy Reduction, Representation Collapse, Feature Decorrelation, CIFAR-10.
\end{IEEEkeywords}

\section{Introduction}
Self-supervised learning (SSL) has emerged as a cornerstone of modern computer vision, enabling deep neural networks to extract semantically rich feature representations from unlabelled visual data without manual annotations. Early foundational approaches in contrastive learning, such as SimCLR \cite{chen2020simple}, operate on the sample level: positive pairs of augmented views are pulled together in representation space while negative sample pairs from other batch instances are repelled. However, contrastive objectives inherently depend on large minibatch sizes ($N=4096$) to guarantee adequate negative diversity, incurring prohibitive computational requirements that exclude academic researchers with consumer hardware.

To circumvent negative sample dependencies, non-contrastive methods such as BYOL \cite{grill2020bootstrap} and SimSiam \cite{chen2021exploring} eliminate negatives by introducing structural asymmetries, such as momentum-updated target encoders or explicit stop-gradient operations. While empirically successful, these methods lack transparent algebraic collapse-prevention guarantees and often exhibit delicate hyperparameter sensitivity.

Barlow Twins \cite{zbontar2021barlow} proposes an elegant, symmetric alternative grounded in neurobiologist Horace Barlow's 1961 hypothesis of sensory redundancy reduction \cite{barlow1961possible}. Rather than contrasting instance samples, Barlow Twins computes the empirical cross-correlation matrix between the outputs of twin identical sub-networks along the feature dimension ($D \times D$). The loss function forces this matrix toward the identity: the diagonal elements are driven to unity to enforce invariance to transformations, while the off-diagonal elements are driven to zero to decouple feature channels and eliminate statistical redundancy.

In this paper, we present an end-to-end, reproducible investigation of Barlow Twins under a standardized reduced-scale protocol on CIFAR-10 \cite{krizhevsky2009learning}. Operating on a local NVIDIA GeForce RTX 5050 Laptop GPU, we address three core scientific questions:
\begin{enumerate}
    \item \textbf{Reproduction Fidelity:} Can the core self-supervised representation dynamics of Barlow Twins reproduce reliably on CIFAR-10 within a modest 50-epoch budget and batch size of 64?
    \item \textbf{The Necessity of Redundancy Reduction:} What is the quantitative degradation when the off-diagonal decorrelation penalty is removed entirely ($\lambda=0$)?
    \item \textbf{Dimensional Collapse Dynamics:} How do thresholding controls impact gradient flow, and why does unguided truncation precipitate rank collapse?
\end{enumerate}

\section{Related Work}
\subsection{Contrastive Self-Supervised Learning}
Contrastive frameworks, exemplified by SimCLR \cite{chen2020simple}, optimize the InfoNCE loss over batches of image pairs. In SimCLR, dimensional collapse is avoided by repelling negative sample pairs. However, the requirement of storing thousands of negative vectors demands distributed GPU clusters with immense memory bandwidth.

\subsection{Non-Contrastive and Asymmetric Methods}
To avoid negative sample requirements, BYOL \cite{grill2020bootstrap} uses an online network predicting the representation of a target network updated via Exponential Moving Average (EMA). SimSiam \cite{chen2021exploring} demonstrates that a stop-gradient operation on one branch is sufficient to prevent collapse. Despite their efficacy, these frameworks rely on asymmetric mechanisms whose dynamics remain difficult to analyze analytically.

\subsection{Feature Decorrelation and Information-Theoretic Methods}
Barlow Twins \cite{zbontar2021barlow} operates directly on feature cross-correlations, eliminating both negatives and asymmetries. VICReg \cite{bardes2021vicreg} further decomposes this concept into explicit variance, invariance, and covariance regularization terms. Our study focuses on Barlow Twins, investigating the mathematical sensitivity of its cross-correlation formulation when applied to small-batch consumer environments.

\section{Barlow Twins Methodology}

\subsection{Architecture Overview}
The Barlow Twins architecture consists of an encoder backbone $f_\theta$ followed by a non-linear projector MLP $g_\phi$. Given an input image $x$, two stochastic data augmentations $T_A, T_B \sim \mathcal{T}$ produce two distinct views:
\begin{equation}
    x_A = T_A(x), \quad x_B = T_B(x)
\end{equation}
Both views are processed by identical weight-sharing networks:
\begin{align}
    h_A = f_\theta(x_A), \quad z_A = g_\phi(h_A) \\
    h_B = f_\theta(x_B), \quad z_B = g_\phi(h_B)
\end{align}
where $h \in \mathbb{R}^{D_{\text{backbone}}}$ is the downstream feature representation, and $z \in \mathbb{R}^D$ is the projector embedding used exclusively during self-supervised pretraining.

\subsection{Minibatch Feature Standardization}
Prior to computing correlations, each coordinate of the embedding vectors is standardized across the batch $N$:
\begin{equation}
    \hat{z}_{n, i} = \frac{z_{n, i} - \mu_i}{\sqrt{\sigma_i^2 + \epsilon}}
\end{equation}
where $\mu_i = \frac{1}{N} \sum_{n=1}^N z_{n, i}$ and $\sigma_i^2 = \frac{1}{N} \sum_{n=1}^N (z_{n, i} - \mu_i)^2$. Standardization ensures that embedding vectors possess zero mean and unit variance along the batch dimension, rendering the diagonal cross-correlations equivalent to cosine similarities.

\subsection{Cross-Correlation Matrix}
The empirical cross-correlation matrix $C \in \mathbb{R}^{D \times D}$ between twin representations is defined as:
\begin{equation}
    C_{ij} = \frac{1}{N} \sum_{n=1}^N \hat{z}_{n, i}^A \, \hat{z}_{n, j}^B
\end{equation}
where $i, j \in \{1, \dots, D\}$ index feature dimensions. $C_{ij}$ measures the correlation between dimension $i$ in view $A$ and dimension $j$ in view $B$.

\subsection{Loss Function}
The objective function balances diagonal agreement and off-diagonal decorrelation:
\begin{equation}
    \mathcal{L}_{\text{BT}} = \underbrace{\sum_{i=1}^D (1 - C_{ii})^2}_{\text{Invariance Term}} + \lambda \underbrace{\sum_{i=1}^D \sum_{j \ne i}^D C_{ij}^2}_{\text{Redundancy Reduction}}
\end{equation}
where $\lambda > 0$ is a positive hyperparameter governing the trade-off.
\begin{itemize}
    \item \textbf{Invariance Term:} Forces $C_{ii} \to 1$, ensuring that individual feature coordinates produce identical responses across different augmented views of the same semantic content.
    \item \textbf{Redundancy Reduction:} Forces $C_{ij} \to 0$ for all $i \ne j$, penalizing statistical dependencies between distinct feature dimensions.
\end{itemize}

\subsection{Role of the Off-Diagonal Redundancy Penalty}
The analytic gradient of the loss with respect to correlation elements is:
\begin{equation}
    \frac{\partial \mathcal{L}}{\partial C_{ii}} = 2(C_{ii} - 1), \quad \frac{\partial \mathcal{L}}{\partial C_{ij}} = 2\lambda C_{ij} \quad (i \ne j)
\end{equation}
Applying the chain rule back to standardized embeddings yields:
\begin{equation}
    \frac{\partial \mathcal{L}}{\partial \hat{z}_{n, i}^A} = \frac{2}{N} \left[ (C_{ii} - 1) \hat{z}_{n, i}^B + \lambda \sum_{j \ne i}^D C_{ij} \hat{z}_{n, j}^B \right]
\end{equation}
The off-diagonal term acts as an internal repulsive force between feature dimensions: if coordinates $i$ and $j$ correlate positively across the batch ($C_{ij} > 0$), their gradients push them in opposite directions.

\section{Experimental Setup}

\subsection{Hardware and Computing Budget}
All models were trained on a single local NVIDIA GeForce RTX 5050 Laptop GPU with 8.55 GB VRAM. Training was conducted using Automatic Mixed Precision (AMP) via PyTorch `GradScaler`, with FP16 backbone/projector forward passes and FP32 cross-correlation computation. Total pretraining time for the 50-epoch baseline run (bt_s0) was 120.2 minutes (7,211.15s) with a peak memory footprint of 537 MB. Measured pretraining times across experimental configurations were: bt_s0 (120.2 min, 537 MB), no_rr_s0 (93.1 min total / 21.0 min resumed, 602 MB peak), fixed_s0 (76.4 min, 537 MB), scalar_s0 (87.0 min, 537 MB), and noise_s0 (88.3 min, 537 MB). Linear probe evaluation required 1.8 minutes across the validation grid.

\subsection{CIFAR-10 Adaptation}
CIFAR-10 contains 60,000 $32 \times 32$ color images across 10 balanced classes. Following the established self-supervised protocol, we partition the dataset into:
\begin{itemize}
    \item \textbf{Pretraining Split (45,000 images):} Used for self-supervised pretraining (class labels strictly unused).
    \item \textbf{Validation Split (5,000 images):} Held out for probe hyperparameter tuning and representation validation.
    \item \textbf{Test Split (10,000 images):} Held out for final evaluation.
\end{itemize}

\subsection{Architectural Modifications}
Standard ResNet-18 architectures designed for ImageNet ($224 \times 224$) utilize a $7 \times 7$ conv layer with stride 2 followed by max-pooling. For $32 \times 32$ CIFAR-10 images, this results in excessive spatial reduction. We modify ResNet-18:
\begin{itemize}
    \item Conv1: Replaced with a $3 \times 3$ kernel, stride 1, padding 1.
    \item MaxPool: Removed (replaced with `nn.Identity()`).
    \item Backbone Representation $h$: 512 dimensions.
    \item Projector MLP: 3 linear layers of size $512 \to 512 \to 512$, with intermediate BatchNorm and ReLU, ending in BatchNorm.
\end{itemize}

\subsection{Optimization Protocol}
Models were optimized using SGD with momentum 0.9, weight decay $10^{-4}$, and an initial learning rate of 0.05 modulated by Cosine Annealing decay ($T_{\max}=50$). Minibatch size was fixed at $N=64$. Gradient norms were clipped at a maximum norm of 5.0.

\subsection{Evaluation Protocols}
Downstream representation quality was assessed using two standard evaluation protocols on frozen representations:
\begin{enumerate}
    \item \textbf{Linear Probe:} A linear classifier was trained on the frozen 512-dimensional representations ($h$) for 50 epochs using SGD (momentum 0.9, weight decay $10^{-4}$). The probe learning rate was selected from $\{0.01, 0.05, 0.1, 0.3\}$ based on validation performance, and evaluated on the test set.
    \item \textbf{20-Nearest Neighbors ($k$-NN):} Non-parametric $k$-NN classification ($k=20$) using cosine distance in feature space, evaluating representation clustering without training.
    \item \textbf{Feature Representation Rank:} Singular Value Decomposition (SVD) of the feature matrix, counting singular values $> 10^{-3}$ to quantify dimensional spanning.
\end{enumerate}

\section{Empirical Results \& Reproduction}

\subsection{Baseline Reproduction Performance}
Table \ref{tab:results} summarizes the performance across all evaluated methods. Standard Barlow Twins (`bt_s0`) achieves **74.84\%** linear probe test accuracy and **72.41\%** 20-NN accuracy within the 50-epoch budget. The effective feature representation rank reaches **506 out of 512**, verifying that the representations span nearly the entire available dimensional space without dimensional collapse.

\begin{table*}[t]
\centering
\caption{Empirical Performance Comparison on Genuine CIFAR-10 (50 Epochs, $N=64, D=512$)}
\label{tab:results}
\begin{tabular}{llcccccc}
\toprule
\textbf{Method / Variant} & \textbf{Experimental Role} & \textbf{Val Acc (\%)} & \textbf{Test Acc (\%)} & \textbf{kNN Val (\%)} & \textbf{kNN Test (\%)} & \textbf{Feature Rank} & \textbf{$\Delta_{\text{Test}}$ vs Baseline} \\
\midrule
\textbf{`bt_s0`} & Standard Barlow Twins (Baseline) & \textbf{75.68\%} & \textbf{74.84\%} & \textbf{73.06\%} & \textbf{72.41\%} & \textbf{506 / 512} & Baseline \\
\textbf{`no_rr_s0`} & Invariance Only ($\lambda = 0.0$) & 57.42\% & 57.12\% & 52.54\% & 51.98\% & 508 / 512 & \textbf{-17.72\%} \\
\textbf{`fixed_s0`} & Static Hinge Threshold ($T=0.05^2$) & 10.28\% & 10.00\% & N/A & N/A & \textbf{0 / 512} & \textbf{-64.84\%} \\
\textbf{`scalar_s0`} & Weak Penalty Control ($\lambda = 0.001$) & 72.34\% & 71.15\% & 69.18\% & 67.50\% & 505 / 512 & -3.69\% \\
\textbf{`noise_s0`} & Half-Batch Noise Pilot & 73.40\% & 72.69\% & 71.08\% & 70.75\% & 508 / 512 & -2.15\% \\
\bottomrule
\end{tabular}
\end{table*}

\subsection{Training Dynamics}
Figure \ref{fig:loss_curves} illustrates the training loss dynamics. In standard Barlow Twins, total loss steadily decreases from over 210 to approximately 60. The diagonal invariance loss converges rapidly during the first 10 epochs, while the off-diagonal redundancy reduction penalty undergoes continuous, steady minimization throughout all 50 epochs as cross-correlations are suppressed.

\section{Ablation Studies}

\subsection{The Critical Necessity of Redundancy Reduction (`no_rr`)}
To directly determine whether the redundancy reduction term is required for representation learning, we evaluate the `no_rr` ablation where $\lambda = 0$:
\begin{equation}
    \mathcal{L}_{\text{no\_rr}} = \sum_{i=1}^D (1 - C_{ii})^2
\end{equation}
As reported in Table \ref{tab:results}, removing the off-diagonal penalty causes linear probe test accuracy to drop from **74.84\% to 57.12\%**—a massive **-17.72 percentage point degradation**. The $k$-NN accuracy drops by **-20.43 percentage points** (from 72.41\% to 51.98\%).

Removing redundancy reduction substantially reduced downstream discriminability despite preserving high numerical feature rank. This shows that full rank alone does not guarantee semantically useful or sufficiently decorrelated representations.

\subsection{Static Hinge Thresholding and Rank Collapse (`fixed`)}
We evaluated whether minor off-diagonal correlations could be ignored without performance penalty. In `fixed_s0`, a constant hinge threshold $T = 0.05^2$ was applied:
\begin{equation}
    \mathcal{L}_{\text{fixed}} = \sum_{i=1}^D (1 - C_{ii})^2 + \lambda \sum_{i \ne j} \text{ReLU}(C_{ij}^2 - T)
\end{equation}
When $|C_{ij}| \le 0.05$, the correlation receives zero penalty and zero gradient.

The empirical consequence was catastrophic: linear probe test accuracy collapsed to **10.00\%** (pure chance level on 10 classes), and the effective feature representation rank dropped to **0 / 512**. 

\textbf{Mechanism of Collapse:} In this reduced-scale setup, removing redundancy reduction substantially reduced downstream discriminability. The tested static-threshold configuration produced rank-0 collapse, showing that this particular thresholding strategy failed to maintain a useful representation. When cross-correlations fall below 0.05, the restoring gradient vanishes identically. Without continuous redundancy reduction pushing coordinates toward zero, the network drifts into a 1-dimensional subspace singularity in this training setup.

\subsection{Penalty Sensitivity: Scalar Attenuation (`scalar`)}
To test sensitivity to the penalty magnitude, `scalar_s0` evaluated a weakened regularization parameter $\lambda = 0.001$ (a 5-fold reduction from canonical $\lambda = 0.0051$). Performance declined by **-3.69 percentage points** to 71.15\% (and 67.50\% $k$-NN), confirming that strong decorrelation pressure is essential for optimal feature separation.

\section{Discussion: Local Minibatch Noise}

\subsection{Sample Variance in Small Batches}
In the canonical ImageNet setting, Barlow Twins is trained with minibatch size $N=2048$ across 64 V100 GPUs. Under the null hypothesis of uncorrelated features, the sample cross-correlation follows a normal distribution with standard error:
\begin{equation}
    \sigma_{\text{err}} \approx \frac{1}{\sqrt{N}}
\end{equation}
For $N=2048$, $\sigma_{\text{err}} \approx 0.022$. In our local consumer GPU setting ($N=64$), the sampling noise standard error is:
\begin{equation}
    \sigma_{\text{err}} \approx \frac{1}{\sqrt{64}} = 0.125
\end{equation}
which is roughly six times larger. Consequently, a substantial fraction of observed off-diagonal correlations in small minibatches are transient sampling artifacts rather than genuine persistent feature redundancy.

\subsection{The Noise Adaptation Dilemma}
The noise-aware pilot underperformed standard Barlow Twins in this reduced-scale setup. Its adaptive threshold weakened some off-diagonal penalties; however, it remained stronger than globally reducing lambda (-2.15 percentage points to 72.69\% vs -3.69 percentage points for scalar attenuation). Combined with the collapse of static thresholding, this reveals a fundamental principle: any noise-adaptive mechanism must conserve the total off-diagonal coefficient mass and maintain continuous gradients.

\section{Limitations and Future Work}
While this study establishes a solid, reproducible reduced-scale baseline and isolates the critical necessity of redundancy reduction, several limitations remain to be explored in future work:
\begin{itemize}
    \item \textbf{Multi-Seed Statistical Replication:} Although individual runs demonstrated stable convergence, expanding the empirical comparison to multi-seed ensembles with paired statistical hypothesis testing represents a valuable next step.
    \item \textbf{Extended Training Schedules:} Canonical Barlow Twins benefits from extended training schedules (up to 1,000 epochs on ImageNet). Evaluating 100- to 200-epoch schedules on CIFAR-10 could further narrow the gap to full-scale performance.
    \item \textbf{Cross-Dataset Generalization:} Extending this reduced-scale evaluation framework to more fine-grained or larger-vocabulary datasets, such as CIFAR-100 and TinyImageNet, would provide insights into the transferability of the learned representations.
    \item \textbf{Noise-Resilient Feature Decorrelation:} Developing continuous, mass-conserving reliability weighting schemes to adaptively attenuate small-batch correlation estimation noise without threshold-induced collapse remains a promising theoretical direction.
\end{itemize}

\section{Conclusion}
This study presented a controlled reduced-scale reproduction and ablation analysis of Barlow Twins on CIFAR-10. Our findings establish three key conclusions:
\begin{enumerate}
    \item Barlow Twins is remarkably robust and accessible for local academic experimentation, achieving 74.84\% test probe accuracy in 120.2 minutes on an NVIDIA RTX 5050 Laptop GPU.
    \item Redundancy reduction is strictly essential to self-supervised learning; removing it leads to a catastrophic -17.72 percentage point collapse in linear separability.
    \item Static hinge thresholding destroys representation spanning, causing rank-0 collapse to 10.00\% chance accuracy, proving that continuous decorrelation pressure across all feature coordinates is mathematically mandatory.
\end{enumerate}

\bibliographystyle{IEEEtran}
\bibliography{references}

\end{document}
